\section{Barrier Synchronisation}

We now implement a barrier synchronisation object using conditions.  Recall
that the object is to be used by~|n| threads.  Each thread calls the operation
|sync|, passing in their own identity (in fact, this implementation won't use
those identities, but we want to give the different implementations of a
barrier the same interface).  Each calls to |sync| should block until all |n|
threads have called it, and then all should return.  

One challenge we face is to stop two different synchronisations interfering:
threads might call |sync| for one synchronisation before all threads have
returned from the previous one.  To this end, we will split each
synchronisation into two phases: \emph{entering} and \emph{leaving}.  If a
thread calls |sync| during the leaving phase of the previous synchronisation,
it waits for that phase to complete.   

%%%%%

\begin{figure}
\begin{scala}
class ConditionsBarrier(n: Int) extends BarrierT{
  /** The number of threads in the main part of the synchronisation. */
  private var count = 0

  /** Are we in the leaving phase of a synchronisation? */
  private var leaving = false

  /** Lock for mutual exclusion. */
  private val lock = new Lock

  /** Condition to indicate that the synchronisation has happened and threads
    * should continue. */
  private val continue = lock.newCondition

  /** Condition to indicate that the previous synchronisation has finished, and
    * threads can enter the main part of the synchronisation. */
  private val enter = lock.newCondition

  def sync(me: Int) = lock.mutex{
    enter.await(!leaving) // Wait for previous synchronisation to end (1).
    if(count == n-1){      // Start leaving phase
      leaving = true; continue.signalAll() // Signal to threads at (2).
    }
    else{      // Have to wait
      count += 1; continue.await() // Wait for the others (2).
      assert(leaving); count -= 1
    }
    if(count == 0){      // Start entering phase.
      leaving = false; enter.signalAll() // Signal to threads at (1).
    }
  }
}
\end{scala}
\caption{A barrier synchronisation implemented using {\scalastyle
    Conditions}.}
\label{fig:ConditionsBarrier}
\end{figure}

%%%%%

The implementation is in Figure~\ref{fig:ConditionsBarrier}.  The variable
|leaving| indicates whether the system is in the leaving phase; the variable
|count| records how many threads are part of the current synchronisation.  We
use a condition |continue| during the entering phase: threads wait on it until
all the threads have called |sync|.  The last (|n|th) thread to enter starts
the leaving phase, and signals to all the threads waiting on |continue|.  We
use a condition |enter| during the leaving phase: threads wait on it until all
threads from the previous synchronisation have returned.  The last thread to
leave starts the entering phase, and signals to all the threads waiting on
|enter|.

\begin{instruction}
Study the details of the implementation. 
\end{instruction}
